\documentclass[11pt]{article}
\usepackage[margin=1in]{geometry}
\usepackage{amsmath,amssymb,amsthm,bm,hyperref,enumitem}
\theoremstyle{plain}
\newtheorem{theorem}{Theorem}\newtheorem{proposition}{Proposition}\newtheorem{corollary}{Corollary}
\newtheorem{lemma}{Lemma}
\theoremstyle{definition}\newtheorem{assumption}{Assumption}\newtheorem{remark}{Remark}
\newcommand{\R}{\mathbb{R}}\newcommand{\E}{\mathbb{E}}\newcommand{\N}{\mathcal{N}}
\newcommand{\Pf}{\mathbb{P}}\newcommand{\I}{\mathbb{1}}\newcommand{\HH}{\mathbb{H}}\newcommand{\TV}{\mathrm{TV}}
\newcommand{\Lpath}{L^\star_{\mathrm{path}}}\newcommand{\Lstr}{L^\star_{\mathrm{struct}}}

\title{\textbf{Structural De-aliasing: Complete Proofs (Part I)}\\[3pt]
\large Full proofs of the results that are ready; rigorous lemmas for the field rate}
\author{Robert Richardson}\date{\today}

\begin{document}\maketitle
\noindent This companion gives complete, self-contained proofs of the results that are ready for full rigor
(\S\ref{sec:rate}--\S\ref{sec:design}), and the rigorous component lemmas for the field-minimax rate
(\S\ref{sec:field}), with the two remaining \emph{constants} (the exact field exponent and the boundary-posterior
scale) precisely delimited. Statements match the submission skeleton.

\section{Setup and assumptions}
Let $\mathcal T=(V,\mathcal E)$ be a tree. Latent states $\{Z_v\}_{v\in V}$, $Z_v\in[K]$, form a Markov random
field with edge transition matrix $T$ ($T_{cc'}=\Pf(Z_u{=}c'\mid Z_v{=}c)$ for $u\sim v$) and an initial law
$\pi$; a measurement channel emits $Y_v\mid Z_v{=}c\sim O_c$ (row $c$ of a $K\times M$ stochastic matrix $O$),
conditionally independent across $v$ given $\{Z_v\}$. Fix an $h$-splitting node-aliased pair $a,b$: $h(a)\neq h(b)$
and $O_a=O_b$ (exactly, unless a node separation $\delta_E>0$ is stated). For $c\in\{a,b\}$ and $\ell\ge1$ write the
\emph{branch signature} $G_c^{(\ell)}=\mathcal L(Y_{\mathrm{sub}}\mid Z_v{=}c)$ for a depth-$\ell$ subtree hanging
off a neighbour of $v$; $G_c^{(1)}=(TO)_c$. Squared Hellinger separations: $\delta_E^2=H^2(O_a,O_b)$,
$\delta_G^2=H^2\!\big(G_a^{(1)},G_b^{(1)}\big)$. For distributions $p,q$ we use the Hellinger affinity
$\rho(p,q)=\sum\sqrt{p\,q}$ (so $H^2=2(1-\rho)$), the Chernoff divergence $C(s)=-\log\sum p^s q^{1-s}$, and the
Chernoff information $C^\star=\max_{s\in[0,1]}C(s)$.

\paragraph{Two problems (kept sharply distinct).} We separate the \emph{oracle} problem, where the channel
$\vartheta=(T,O)$ is \emph{known} and only the latent configuration $\{Z_v\}$ is unknown, with risk
$\mathcal R_{\mathrm{oracle}}(\vartheta)=\inf_{\hat R}\E_\vartheta\,\mathcal L(\hat R,R)$; from the \emph{minimax}
problem, where the channel itself is unknown, $\mathcal R_n(\Theta)=\inf_{\hat R}\sup_{\vartheta\in\Theta}
\E_\vartheta\,\mathcal L(\hat R,R)$. Bayes classifiers using the exact $(P_a,P_b)$ are available only in the oracle
problem; the gap between the two is the singular estimation remainder of \S\ref{sec:singular}.
\begin{lemma}[Marginal-mode optimality]\label{thm:oracle-risk}
For additive Hamming loss the Bayes rule is coordinatewise, $\hat R_v=\arg\max_r\Pf(R_v{=}r\mid Y)$; on a tree
these marginals are computed exactly by belief propagation, and $\mathcal R_{\mathrm{oracle}}(\vartheta)=\frac1n
\sum_v P_e(v)$ with $P_e(v)$ the nodewise two-point Bayes error.
\end{lemma}
\begin{proof}
Additivity of the loss makes the risk $\frac1n\sum_v\Pf(\hat R_v\neq R_v)$ separately minimized by the marginal
posterior mode at each $v$; exactness of belief propagation on a tree is standard.
\end{proof}

\section{The near-aliasing rate}\label{sec:rate}
Consider testing $H_a:Z_v{=}a$ against $H_b:Z_v{=}b$ (equal prior) from the node measurement $Y_v$ and $d$
neighbour views at radius $L$; write $G_c=G_c^{(L)}$, with $C_E(s),C_G(s)$ the Chernoff divergences of
$(O_a,O_b)$ and $(G_a,G_b)$, and $\rho_E=e^{-C_E(1/2)},\rho_G=e^{-C_G(1/2)}$.

\begin{lemma}[Conditional independence / product factorization]\label{lem:prod}
Given $Z_v=c$, the node emission $Y_v$ and the observations in the $d$ distinct branches at $v$ are mutually
independent, and each branch has law $G_c$. Hence under $H_c$ the joint observation law is
$P_c=O_c\otimes G_c^{\otimes d}$, and for every $s$,
$\int P_a^s P_b^{1-s}=e^{-C_E(s)}\,e^{-d\,C_G(s)}=e^{-(C_E(s)+d\,C_G(s))}$.
\end{lemma}
\begin{proof}
Deleting $v$ from the tree disconnects the branches and the singleton $\{v\}$; by the global Markov property of a
tree-indexed MRF, the observations in different branches and $Y_v$ are mutually independent given $Z_v$. Each
branch subtree, rooted at a neighbour whose state is drawn from row $c$ of $T$, has observation law $G_c$ by
definition, and the $d$ branches are i.i.d.\ copies given $Z_v=c$. The $s$-tilted affinity of a product measure
factorizes, $\int(\prod_i p_i)^s(\prod_i q_i)^{1-s}=\prod_i\int p_i^s q_i^{1-s}$, giving the stated product.
\end{proof}

\begin{theorem}[Near-aliasing rate]\label{thm:rate}
The equal-prior Bayes error $P_e$ satisfies
\begin{equation}\label{eq:twoside}
   \tfrac14\big(\rho_E\rho_G^{d}\big)^2\ \le\ P_e\ \le\ \tfrac12\,\rho_E\rho_G^{d}.
\end{equation}
The \emph{exact} consistency condition is $-\log\rho_E-d\log\rho_G\to\infty$; \emph{near the alias}, where
$-\log\rho_\bullet=\tfrac12\delta_\bullet^2(1+o(1))$, this becomes $\delta_E^2+d\,\delta_G^2\to\infty$ and
$P_e=\exp(-\Theta(\delta_E^2+d\delta_G^2))$. The sharp error exponent is the product Chernoff information with a
\emph{shared} tilt,
\begin{equation}\label{eq:sharp}
   -\log P_e = \max_{s\in[0,1]}\{C_E(s)+d\,C_G(s)\}\,(1+o(1)),
\end{equation}
and near the alias $s^\star\to\tfrac12$ so the exponent is $\tfrac12(\delta_E^2+d\,\delta_G^2)(1+o(1))$.
\end{theorem}
\begin{proof}
Write $\rho=\rho(P_a,P_b)=\int\sqrt{P_aP_b}=\rho_E\rho_G^{d}$ by Lemma~\ref{lem:prod} at $s=\tfrac12$. The
equal-prior Bayes error is $P_e=\tfrac12\int\min(P_a,P_b)$.
\emph{Upper bound:} $\min(p,q)\le\sqrt{pq}$ pointwise, so $P_e\le\tfrac12\int\sqrt{P_aP_b}=\tfrac12\rho$.
\emph{Lower bound:} let $M=\int\min(P_a,P_b)$ and $\int\max=2-M$. Cauchy--Schwarz gives
$\rho=\int\sqrt{\min}\sqrt{\max}\le\sqrt{M}\sqrt{2-M}\le\sqrt{2M}$, so $M\ge\rho^2/2$ and $P_e\ge\rho^2/4$. This is
\eqref{eq:twoside}. Since $-\log\rho_\bullet=-\log(1-\tfrac12\delta_\bullet^2)=\tfrac12\delta_\bullet^2(1+o(1))$,
$-\log\rho=\tfrac12(\delta_E^2+d\delta_G^2)(1+o(1))$ and both sides of \eqref{eq:twoside} are
$\exp(-\Theta(\delta_E^2+d\delta_G^2))$.
\emph{Sharp exponent:} the log-likelihood ratio $\Lambda=\log\frac{dP_b}{dP_a}$ is, by Lemma~\ref{lem:prod}, a sum
of the node term and $d$ i.i.d.\ branch terms. By Chernoff's theorem the equal-prior Bayes-error exponent equals
the Chernoff information of the experiment, $\max_s-\log\int P_a^sP_b^{1-s}=\max_s\{C_E(s)+dC_G(s)\}$ (one $s$
tilts the whole product), giving \eqref{eq:sharp}. Near the alias each $C_\bullet(s)$ is maximized at $s=\tfrac12$
to leading order with $C_\bullet^\star=\tfrac12\delta_\bullet^2(1+o(1))$, giving the stated exponent.
\end{proof}
\begin{remark}[The polynomial prefactor]
For a \emph{fixed} branch channel as $d\to\infty$, Bahadur--Rao gives $P_e\asymp\{d\,V(s^\star)\}^{-1/2}
\exp(-\max_s\{C_E(s)+dC_G(s)\})$ under a nonlattice condition, where $V(s^\star)$ is the tilted log-likelihood-ratio
variance; this is the $d^{-1/2}$ form. If instead the branch channel itself approaches aliasing with $d$, then
$V\to0$ and the prefactor is more accurately $\{d\,V_d(s_d^\star)\}^{-1/2}$, typically of order
$(d\,\delta_G^2)^{-1/2}$. The \emph{exponent} \eqref{eq:sharp} is robust across both regimes; only the prefactor
changes.
\end{remark}

\begin{proposition}[Branching de-aliasing]\label{prop:branch}
There is a finite two-layer latent model with node-aliased $a,b$ in which \emph{every constituent single
measurement is identically distributed} under both classes, yet the two-sibling \emph{joint} differs. Thus the
joint structural law separates states that no individual measurement does; the de-aliasing exponent is the positive
Chernoff information of the joint signatures, while the per-measurement Chernoff information is $0$.
\end{proposition}
\begin{proof}
Binary measurement. A centre class $c\in\{a,b\}$ acts through a hidden intermediate $U\in\{1,2,3\}$ with
$\Pf(U{=}j\mid c)=T_c(j)$; conditional on $U{=}j$ the observed children are i.i.d.\ $\mathrm{Bernoulli}(m_j)$. Take
$m=(0.2,0.5,0.8)$, $T_a=(\tfrac12,0,\tfrac12)$, $T_b=(0,1,0)$. A single child has mean
$\bar m_c=\sum_j T_c(j)m_j$: $\bar m_a=\tfrac12(0.2)+\tfrac12(0.8)=0.5=\bar m_b$, so each single observed
measurement is $\mathrm{Bernoulli}(0.5)$ under both classes and the per-measurement Chernoff information is $0$. The
pair correlation is $\E[Y_1Y_2\mid c]=\sum_jT_c(j)m_j^2$: $0.34$ for $a$, $0.25$ for $b$, so the two-sibling joints
differ and their Chernoff information is strictly positive (the de~Finetti mixtures $\sum_jT_c(j)\,
\mathrm{Bern}(m_j)^{\otimes D}$ differ).
\end{proof}
\begin{remark}
This proves ``all constituent marginals identical, joint different'', hence $\Lstr<\infty$ with the per-measurement
signature aliased. It does \emph{not} assert $(T^\ell O)_a=(T^\ell O)_b$ for all $\ell$ in a homogeneous recurrent
chain (which would need a trace-equivalent, non-bisimilar pair from the realization-theory literature); as the
linear-vs-branching gap is scaffolding here, we state only the finite two-layer phenomenon.
\end{remark}

\section{The information limit}\label{sec:info}
Let $X_v$ be covariates and $\mathcal I_G=I\!\big(R_v;Y_{N(v)}\mid Y_v,X_v\big)$.
\begin{theorem}[Information limit; quantitative]\label{thm:info}
Under log loss the reduction in optimal risk from adding $Y_{N(v)}$ equals $\mathcal I_G$. For \emph{any} loss
bounded by $1$, the reduction in optimal Bayes risk is at most $\sqrt{\mathcal I_G/2}$; in particular
$\mathcal I_G=0$ implies no neighbourhood rule improves optimal prediction of $R_v$.
\end{theorem}
\begin{proof}
The log-loss optimal risk given a $\sigma$-field is the corresponding conditional entropy, so the reduction is
$\HH(R_v\mid Y_v,X_v)-\HH(R_v\mid Y_v,X_v,Y_{N(v)})=\mathcal I_G$ by definition of conditional mutual information.
For a bounded loss, the Bayes-risk change from replacing the posterior $\Pf(R_v\mid Y_v,X_v)$ by
$\Pf(R_v\mid Y_v,X_v,Y_{N(v)})$ is at most $\E\big\|\Pf(R_v\mid Y_v,X_v,Y_{N(v)})-\Pf(R_v\mid Y_v,X_v)\big\|_\TV$.
By the conditional Pinsker inequality and Jensen (concavity of $\sqrt{\cdot}$),
\[
   \E\|\cdot\|_\TV\ \le\ \E\sqrt{\tfrac12\mathrm{KL}\big(\Pf(R_v\mid Y_v,X_v,Y_N)\,\|\,\Pf(R_v\mid Y_v,X_v)\big)}
   \ \le\ \sqrt{\tfrac12\,\E[\mathrm{KL}]}\ =\ \sqrt{\mathcal I_G/2},
\]
using $\E[\mathrm{KL}]=\mathcal I_G$.
\end{proof}

\section{Functional identifiability}\label{sec:func}
\begin{theorem}[Functional identifiability, oriented model]\label{thm:func}
Fix an oriented parameterization (a labelling of the states). Then the \emph{model-identifiable} functionals
$R=h(Z)$ recoverable from the radius-$L$ neighbourhood law are exactly those constant on the depth-$L$
observational-equivalence classes: $h$ is identified iff every $h$-splitting pair $c,c'$ has $Q_L(c)\neq Q_L(c')$.
Such an $h$ can be identified even when $Z$ and $O$ are not.
\end{theorem}
\begin{proof}
$h$ is identified iff the map $Q_L(c)\mapsto h(c)$ is well defined, i.e.\ $Q_L(c)=Q_L(c')\Rightarrow h(c)=h(c')$;
the contrapositive is that every $h$-splitting pair is de-aliased. As $L\to\infty$, $Q_\infty(c)=Q_\infty(c')$ is
observational equivalence $=$ the coarsest emission-consistent lumpable partition (bisimulation$=$lumpability), and
$h$ is identified iff measurable with respect to it. This constrains only $h$ relative to the partition and is
compatible with $Z$ (finer) and $O$ being non-identified.
\end{proof}
\begin{remark}[Orientation and per-draw recovery]
Two caveats. \emph{(i) Unoriented models.} The statement fixes a labelling. In an unoriented latent model a global
label automorphism preserving the observed law can change $h$ (the sign/orientation of \S\ref{sec:singular});
observational identification then additionally requires $h$ to be invariant under all such automorphisms, or an
anchor to select the orientation. \emph{(ii) Model vs.\ realization.} Distinct $Q_L(c),Q_L(c')$ make the latent
classes \emph{statistically distinguishable} (model/functional identifiability); they do not entail that a single
noisy neighbourhood realization determines $R_v$ with probability one---that is the classification problem of
\S\ref{sec:rate}, with error $P_e>0$.
\end{remark}

\section{Forest risk: the exact Chernoff exponent}\label{sec:forest}
Let the field be the centre states of $n$ disjoint depth-$1$ star gadgets ($h(a)\neq h(b)$, so the centre state
determines $R_i$), loss $\mathcal L(\hat R,R)=\frac1n\sum_i\I\{\hat R_i\neq R_i\}$; let $P_e$ be the equal-prior
two-point Bayes error of one gadget.
\begin{theorem}[Forest risk, exact exponent]\label{thm:forest}
With $\vartheta$ known, $P_e\le\mathcal R_{\mathrm{oracle}}(\mathrm{forest})\le 2P_e$; hence the forest risk has
\emph{exactly} the single-gadget Chernoff exponent $\max_s\{C_E(s)+dC_G(s)\}$ (Theorem~\ref{thm:rate}), not merely
an $\exp(-\Theta)$ rate.
\end{theorem}
\begin{proof}
\emph{Lower:} put the uniform prior on $\omega\in\{a,b\}^n$; since the gadgets are independent the Bayes risk is
the average of per-gadget Bayes risks, $=P_e$, and the minimax risk dominates any Bayes risk, so
$\mathcal R\ge P_e$. \emph{Upper:} apply the equal-prior likelihood-ratio test independently on each gadget; if its
two type errors are $(\alpha,\beta)$ then $P_e=\tfrac12(\alpha+\beta)$ and the per-gadget error is at most
$\max(\alpha,\beta)\le\alpha+\beta=2P_e$, uniformly in $\omega$, so $\mathcal R\le 2P_e$. Both bounds carry the
exponent of $P_e$, which is the single-gadget Chernoff information by Theorem~\ref{thm:rate}.
\end{proof}
\begin{remark}[Known vs.\ unknown channel]
The upper bound uses the likelihood-ratio test built from the true $(P_a,P_b)$, so Theorem~\ref{thm:forest} is an
\emph{oracle} ($\vartheta$-known) statement. When $\vartheta$ is unknown and estimated from the same forest, the
minimax risk is $\mathcal R_{\mathrm{oracle}}$ plus the singular estimation remainder governed by
$\lambda=\sqrt n\,\delta_G^2$ (\S\ref{sec:singular}).
\end{remark}

\section{The singular local structure}\label{sec:singular}
Parametrize the separation by a scalar $\eta$ with $G_a=\bar G-\eta u$, $G_b=\bar G+\eta u$ ($\sum u=0$),
$O_a=O_b$, so $\eta=0$ is full alias and $\delta_G\propto|\eta|$ near $0$. Assume the equal (label-symmetric)
mixture over the aliased pair and smoothness of $L$ in $\eta$.
\begin{proposition}[Singular structure]\label{prop:singular}
(a) The marginal likelihood is even: $L(\eta)=L(-\eta)$; hence the score at $\eta=0$ vanishes and the Fisher
information for $\eta$ is degenerate. (b) Under the genericity condition $\partial_\psi \tilde L\,|_{\psi=0}\neq0$
(where $\tilde L(\psi)=L(\sqrt\psi)$, $\psi=\eta^2$), $\psi$ is regular with positive information, so
$\sqrt n(\hat\psi-\psi_0)\rightsquigarrow\N(0,I_\psi^{-1})$; consequently at $\psi_0=0$, $\hat\psi=O_p(n^{-1/2})$ and
$\hat\eta=O_p(n^{-1/4})$, and the invariant local index is $\lambda=\sqrt n\,\psi\ (\propto\sqrt n\,\delta_G^2)$.
(c) The sign of $\eta$ (the orientation of the pair) is not identified by the data.
\end{proposition}
\begin{proof}
\emph{(a)} Consider the involution $\sigma$ that swaps the labels $a\leftrightarrow b$ of the aliased pair. Under
$\sigma$, $(G_a,G_b)\mapsto(G_b,G_a)$, i.e.\ $\eta\mapsto-\eta$, while $O_a=O_b$ is fixed; because the prior over
$\{a,b\}$ is symmetric, the map $(\eta,\{Z_v\})\mapsto(-\eta,\sigma\{Z_v\})$ preserves the joint law of the data.
Marginalizing the latent states, which $\sigma$ permutes measure-preservingly, gives $L(\eta)=L(-\eta)$. An even
smooth function has $L'(0)=0$, so the score $\partial_\eta\log L|_0=0$ and $I_\eta(0)=\E[(\partial_\eta\log L)^2]|_0
=0$. \emph{(b)} Evenness lets us write $L(\eta)=\tilde L(\eta^2)=\tilde L(\psi)$. By genericity
$\partial_\psi\tilde L|_0\neq0$, so $\psi$ has a nondegenerate score at the boundary and finite positive
information $I_\psi(0)$; the MLE/posterior for $\psi$ is regular at rate $\sqrt n$. Hence at the alias
$\psi_0=0$, $\hat\psi=O_p(n^{-1/2})$ (truncated at $0$), and $\hat\eta=\sqrt{\hat\psi}=O_p(n^{-1/4})$. Writing
$\psi\propto\delta_G^2$, the local reparametrization $\psi_0=\lambda/\sqrt n$ keeps $\sqrt n\,\psi\to\lambda$ fixed,
i.e.\ $\delta_G\sim n^{-1/4}$. \emph{(c)} By (a) the likelihood assigns equal value to $\eta$ and $-\eta$, i.e.\ to
the two orientations; no function of the data distinguishes them, so $\mathrm{sign}(\eta)$ is not identified and
its posterior equals its prior unless external (anchor) information breaks the symmetry.
\end{proof}
\begin{remark}[Two stages]
Since $h(a)\neq h(b)$, orientation \emph{is} the target: structure identifies $|\eta|,\psi$ and the unordered pair;
a single anchor supplies $\mathrm{sign}(\eta)$ and hence $R=h(Z)$. Orientation therefore costs $O(1)$ external
information regardless of $n$.
\end{remark}

\begin{proposition}[Boundary posterior; independent-star experiment]\label{prop:boundary}
Assume the quadratic-mean expansion $p_\eta(y)=p_0(y)+\eta^2a(y)+o(\eta^2)$ with $\sum_y a(y)=0$ and
$I_\psi=\sum_y a(y)^2/p_0(y)\in(0,\infty)$, and write $\psi=\eta^2\ge0$, $S(y)=a(y)/p_0(y)$. Then for the local
parameter $h=\sqrt n\,\psi$,
\[
   \ell_n(h/\sqrt n)-\ell_n(0)=h\,\Delta_n-\tfrac12 h^2 I_\psi+o_P(1),\qquad
   \Delta_n=\tfrac1{\sqrt n}\sum_{i=1}^n S(Y_i)\Rightarrow \N(0,I_\psi),
\]
and under $\psi_n=\lambda_0/\sqrt n$, Le~Cam's third lemma gives $\Delta_n\Rightarrow\N(\lambda_0 I_\psi,I_\psi)$.
If the prior density of $\psi$ is continuous and positive at $0$, the posterior for $h\ge0$ converges to
\[
   \pi(h\mid\Delta)\ \propto\ \exp\big\{h\Delta-\tfrac12 I_\psi h^2\big\}\,\I\{h\ge0\}
   \ =\ \N\!\big(\Delta/I_\psi,\ I_\psi^{-1}\big)\ \text{truncated to}\ [0,\infty).
\]
\end{proposition}
\begin{proof}
Writing $\psi$ for the regular parameter (Prop.~\ref{prop:singular}(b)) and expanding the log-likelihood in $\psi$
by the differentiability-in-quadratic-mean condition gives the displayed LAN expansion with score $S$ and
information $I_\psi$; $\Delta_n\Rightarrow\N(0,I_\psi)$ by the CLT. Under the contiguous local alternative
$\psi_n=\lambda_0/\sqrt n$, Le~Cam's third lemma shifts the limit mean to $\lambda_0 I_\psi$. With a prior
continuous and positive at $\psi=0$ the prior is locally flat in $h$, so the posterior is proportional to the local
likelihood restricted to the boundary $h\ge0$, giving the stated truncated Gaussian (a boundary Bernstein--von
Mises statement).
\end{proof}
\begin{remark}[Corrected interpretation: magnitude is data-driven, only the sign is not]
At $\lambda_0=0$ the posterior for the \emph{magnitude} $h=\sqrt n\,\psi$ is $\N(\Delta/I_\psi,I_\psi^{-1})$
truncated to $[0,\infty)$ with $\Delta\Rightarrow\N(0,I_\psi)$---a \emph{data-driven} truncated Gaussian, \emph{not}
the prior. So the earlier ``interpolates between Bernstein--von Mises and prior'' description is wrong for the
magnitude: $\psi=\eta^2$ is locally regular at the boundary and data-driven for every $\lambda_0$. The permanently
prior/anchor-sensitive object is the \emph{sign}. Consequently $|\eta|$ lives on the $n^{-1/4}$ scale; the signed
$\eta$ is not estimable at any rate without orientation, and after a perfect anchor it inherits the $n^{-1/4}$
boundary rate.
\end{remark}

\section{The measurement-design law}\label{sec:design}
\begin{corollary}[First-order design law]\label{cor:design}
A design with $n$ replicated node measurements, $d$ radius-$L$ views, and $m$ readings of a second instrument
(Chernoff divergence $C_2$) has error exponent $\max_s\{n\,C_E(s)+d\,C_G(s)+m\,C_2(s)\}$, whose near-alias
first-order value is $n\,C_E+d\,C_G+m\,C_2$; iso-information surfaces give the exact substitution rates.
\end{corollary}
\begin{proof}
The three sources are independent channels, so by Lemma~\ref{lem:prod} the $s$-tilted affinity of the full
experiment is $\exp(-[nC_E(s)+dC_G(s)+mC_2(s)])$; the Bayes-error exponent is the maximum over the shared $s$
(Theorem~\ref{thm:rate}). Near the alias every $C_\bullet(s)$ is maximized at $s=\tfrac12$ to leading order with
$C_\bullet=\tfrac12\delta_\bullet^2(1+o(1))$, giving the linear law.
\end{proof}

\section{The field rate (known channel) via the branch-information sandwich}\label{sec:field}
For \emph{known} $\vartheta=(T,O)$, under average Hamming loss the Bayes rule is coordinatewise marginal posterior
mode (Theorem~\ref{thm:oracle-risk}), so the oracle field risk $\mathcal R_{\mathrm{oracle}}(\vartheta)=\frac1n
\sum_v P_e(v)$ is the average of the nodewise two-point Bayes errors, each using \emph{all} of the tree's
information about $Z_v$. We bound $P_e(v)$ directly by data processing---no spectral gap, packing, or Assouad.
This replaces the spectral-gap saturation argument, which is invalid (contraction along a path does not control
aggregate information on a branching tree, and the branching example of Prop.~\ref{prop:branch} shows
$C_G^{(\infty)}\gg C_G^{(1)}$ is possible when the immediate channel hides information).

For a branch $j$ at $v$, let $U_j$ be the first-neighbour state ($U_j\sim T_c$ given $Z_v{=}c$) and $W_j$ the
\emph{entire} branch of observations, so $P(W_j\mid Z_v{=}c)=(T_c K_j)$ for the branch channel $K_j$; write
$G_c^{(j)}=P(W_j\mid Z_v{=}c)$ and $G_c^{(\infty)}$ its infinite-depth limit.
\begin{lemma}[Branch-information sandwich]\label{lem:sandwich}
For every tilt $s$, $\ C_s\big((TO)_a,(TO)_b\big)\ \le\ C_s\big(G_a^{(j)},G_b^{(j)}\big)\ \le\ C_s(T_a,T_b)$.
\end{lemma}
\begin{proof}
Both inequalities are data processing for the Chernoff divergence, which is monotone under Markov channels.
\emph{Upper:} $G_c^{(j)}=T_cK_j$ is the pair $(T_a,T_b)$ pushed through the channel $K_j$, so
$C_s(T_aK_j,T_bK_j)\le C_s(T_a,T_b)$: the whole branch reveals $Z_v$ at most as well as observing the first hidden
neighbour's state $U_j$ perfectly. \emph{Lower:} the first neighbour's own emission $Y_{U_j}$, with law $(T_cO)$,
is a coarsening (sub-observation) of $W_j$, so $C_s((TO)_a,(TO)_b)\le C_s(G_a^{(j)},G_b^{(j)})$.
\end{proof}
\begin{assumption}[Local observability]\label{ass:obs}
$C_s(T_a,T_b)\le K\,C_s\big((TO)_a,(TO)_b\big)$ for a constant $K<\infty$, uniformly over $s$ and the local model
class. (Equivalently, the immediate channel does not hide the transition-row contrast; $K=1.94$ in our running
model. Its failure is exactly the branching-de-aliasing regime of Prop.~\ref{prop:branch}.)
\end{assumption}
\begin{theorem}[Known-channel field rate]\label{thm:field-rate}
On \emph{any} bounded-degree tree with $d$ branches per node, under Assumption~\ref{ass:obs},
$-\log P_e(v)=\Theta(\delta_E^2+d\,\delta_G^2)$ for every node, hence
$\mathcal R_{\mathrm{oracle}}(\vartheta)\asymp\Phi(\delta_E^2+d\,\delta_G^2)$. No mixing, packing, or Assouad
argument is used.
\end{theorem}
\begin{proof}
By Lemma~\ref{lem:prod} the per-node experiment is the product of the node channel and the $d$ branches, with
$s$-tilted exponent $C_E(s)+\sum_{j=1}^d C_s(G_a^{(j)},G_b^{(j)})$; Theorem~\ref{thm:rate} gives $-\log P_e(v)
\asymp\max_s\{C_E(s)+\sum_j C_s(G_a^{(j)},G_b^{(j)})\}$. Lemma~\ref{lem:sandwich} sandwiches each branch term,
so the whole exponent lies between $\max_s\{C_E(s)+d\,C_s((TO)_a,(TO)_b)\}$ and $\max_s\{C_E(s)+d\,C_s(T_a,T_b)\}$.
Near the alias $C_E^\star\asymp\delta_E^2$ and $C^\star((TO)_a,(TO)_b)\asymp\delta_G^2$, and by
Assumption~\ref{ass:obs} $C^\star(T_a,T_b)\le K\,C^\star((TO)_a,(TO)_b)\asymp\delta_G^2$; both fences are therefore
$\Theta(\delta_E^2+d\,\delta_G^2)$, and so is $-\log P_e(v)$.
\end{proof}
\begin{remark}[Why this is the right theorem]
The sandwich is exact and assumption-light: an infinite branch can reveal $Z_v$ at most as well as its first hidden
neighbour ($U_j$), so deeper structure changes only the constant \emph{unless} the immediate channel is
uninformative relative to the transition rows---i.e.\ unless Assumption~\ref{ass:obs} fails, which is precisely the
branching-de-aliasing regime where deeper observations carry qualitatively new information. The known-channel
field rate thus needs neither a spectral gap nor an antichain; the unknown-channel minimax uses the independent
forest (\S\ref{sec:forest}), where Assouad is exact.
\end{remark}
\begin{remark}[The one remaining constant]
The exact boundary-posterior scale $I_\psi(0)^{-1/2}$ (Prop.~\ref{prop:boundary}) is proved below via LAN for the
independent-star experiment; only its extension to the general tree, and the exact within-$\Theta$ field constant,
remain---both constants, not existence.
\end{remark}

\end{document}
